\section{区间上的Haar展开}
\label{sec:wavelet-interval-expansion}

有限层级的空间$V_J$只含在第$J$层二进区间上取常值的函数。$f(x)=x$这样简单的函数也不在任何$V_J$中，因为它在每个区间内都在变化。要用Haar函数表示一般的$L^2([0,1])$函数，必须让层级无限增加，并说明这样做确实能逼近任意元素。本节先证明尺度空间的并是稠密的，再用第~\ref{chap:inner-product-spaces}章的正交展开得到系数、重构与能量等式。

\subsection{尺度空间的稠密性}
\label{sec:wavelet-scaling-density}

我们要证明：任给$f\in L^2([0,1])$和$\varepsilon>0$，存在某个$j$和$g\in V_j$，使$\lVert f-g\rVert_2<\varepsilon$。$V_j$中的元素称为二进阶梯函数。思路是把一般函数逐步化简：先截去大值，再换成简单函数，再换成有限个区间的指示函数，最后把区间端点移到二进点上。每一步的误差都可以任意小。

\begin{theorem}[二进阶梯函数的稠密性]
\label{thm:wavelet-dyadic-dense}
$\bigcup_{j\geq0}V_j$在$L^2([0,1])$中稠密。因此对每个$f\in L^2([0,1])$，
\[
\lVert f-P_jf\rVert_2\to0\qquad(j\to\infty),
\]
并且$\lVert f-P_jf\rVert_2$随$j$单调不增。
\end{theorem}

\begin{proof}
\emph{第一步：截去大值。}对$M>0$，令$f_M=f\mathbf{1}_{\{\lvert f\rvert\leq M\}}$。逐点有$f_M\to f$（$M\to\infty$），且$\lvert f-f_M\rvert^2\leq\lvert f\rvert^2$可积。由控制收敛定理，$\lVert f-f_M\rVert_2\to0$。所以只需逼近有界可测函数。

\emph{第二步：换成简单函数。}由第~\ref{chap:measure}章的简单函数逼近，有界可测函数$f_M$（复值时分别处理实部与虚部）可以被只取有限个值的简单函数一致逼近。区间$[0,1]$的测度为$1$，一致误差不超过$\delta$时，$L^2$误差也不超过$\delta$。所以只需逼近简单函数，又由线性只需逼近单个指示函数$\mathbf{1}_E$，$E\subset[0,1]$可测。

\emph{第三步：换成有限个区间。}由第~\ref{chap:smoothing-weak-derivatives}章所用的Lebesgue测度的开集外逼近，对$\delta>0$存在开集$U\supset E$，$\lambda(U\setminus E)<\delta$。实轴上的开集是可数个互不相交开区间的并，$U=\bigcup_iJ_i$，且$\sum_i\lambda(J_i)=\lambda(U)<\infty$；取足够多项使剩余部分的测度小于$\delta$，记前$n$个区间之并为$G$。于是$\lambda(E\mathbin{\triangle}G)\leq\lambda(U\setminus E)+\lambda(U\setminus G)<2\delta$。把$G$换成$G\cap[0,1]$不增大对称差，因为$E\subset[0,1]$。指示函数之差只取$0,\pm1$，所以$\lVert\mathbf{1}_E-\mathbf{1}_{G\cap[0,1]}\rVert_2^2=\lambda(E\mathbin{\triangle}(G\cap[0,1]))<2\delta$。

\emph{第四步：端点移到二进点。}对区间$[a,b]\subset[0,1]$，把$a$、$b$分别换成不超过它们的最大的$2^{-j}$整数倍$a_j$、$b_j$。$[a_j,b_j)$是若干个第$j$层二进区间之并，其指示函数属于$V_j$，并且$\lambda([a,b]\mathbin{\triangle}[a_j,b_j))\leq\lvert a-a_j\rvert+\lvert b-b_j\rvert\leq2^{1-j}$，所以$L^2$误差不超过$2^{(1-j)/2}$。有限个区间的指示函数之和，可以对足够大的同一个$j$同时这样处理。四步合起来，任意$f$都能被某个$V_j$中的函数以任意精度逼近。

\emph{收敛与单调性。}给定$\varepsilon>0$，取$j_1$和$g\in V_{j_1}$使$\lVert f-g\rVert_2<\varepsilon$。对$j\geq j_1$，$g\in V_{j_1}\subset V_j$，而$P_jf$是$V_j$中离$f$最近的元素，故$\lVert f-P_jf\rVert_2\leq\lVert f-g\rVert_2<\varepsilon$。同样由$V_j\subset V_{j+1}$得$\lVert f-P_{j+1}f\rVert_2\leq\lVert f-P_jf\rVert_2$。
\end{proof}

证明的最后一段只用到两件事：空间是嵌套的，投影给出最佳逼近。所以“并稠密”自动转化为“投影列收敛”，不需要对每个$f$单独估计$P_jf$。

\begin{example}[线性函数的投影误差]
\label{ex:wavelet-linear-projection-error}
取$f(x)=x$。在每个第$j$层区间上，$P_jf$等于区间中点$m$，记$h=2^{-j}$为区间长度，则
\[
\lVert f-P_jf\rVert_2^2=\sum_{k=0}^{2^j-1}\int_{I_{j,k}}(x-m)^2\,dx=2^j\cdot\frac{h^3}{12}=\frac{4^{-j}}{12},
\]
即$\lVert f-P_jf\rVert_2=2^{-j}/(2\sqrt3)$。层级每增加一层，误差减半。定理~\ref{thm:wavelet-dyadic-dense}只保证误差趋于零，并不给出这样的速度；速度来自$f$本身的光滑性，一般理论见第~\ref{chap:function-approximation}章。
\end{example}

\subsection{系数与重构}
\label{sec:wavelet-coefficients-reconstruction}

稠密性把有限层的分解~\eqref{eq:wavelet-finite-decomposition}推广到整个$L^2([0,1])$。为了与后面实轴上的情形对照，起始层不必是第$0$层：固定整数$j_0\geq0$，由定理~\ref{thm:wavelet-detail-basis}反复分解，对$J>j_0$有$V_J=V_{j_0}\oplus W_{j_0}\oplus\cdots\oplus W_{J-1}$。

\begin{theorem}[区间上的Haar展开]
\label{thm:wavelet-haar-expansion}
固定整数$j_0\geq0$。函数族
\[
\{\varphi_{j_0,k}:0\leq k<2^{j_0}\}\cup\{\psi_{j,k}:j\geq j_0,\ 0\leq k<2^j\}
\]
是$L^2([0,1])$的标准正交基。对每个$f\in L^2([0,1])$，在$L^2$范数下
\begin{equation}
f=\sum_{k=0}^{2^{j_0}-1}\langle\varphi_{j_0,k},f\rangle\varphi_{j_0,k}+\sum_{j=j_0}^{\infty}\sum_{k=0}^{2^j-1}\langle\psi_{j,k},f\rangle\psi_{j,k},
\label{eq:wavelet-haar-expansion}
\end{equation}
并且
\begin{equation}
\lVert f\rVert_2^2=\sum_{k=0}^{2^{j_0}-1}\lvert\langle\varphi_{j_0,k},f\rangle\rvert^2+\sum_{j=j_0}^{\infty}\sum_{k=0}^{2^j-1}\lvert\langle\psi_{j,k},f\rangle\rvert^2 .
\label{eq:wavelet-haar-parseval}
\end{equation}
把$j$求和截到$j\leq J-1$所得的部分和等于$P_Jf$。
\end{theorem}

\begin{proof}
不同层的细节空间$W_j$彼此正交，并都与$V_{j_0}$正交：若$j<j'$，则$W_j\subset V_{j+1}\subset V_{j'}$，而$W_{j'}\perp V_{j'}$。同一层内的正交性由定理~\ref{thm:wavelet-detail-basis}给出，所以整个函数族是标准正交的。它的有限线性组合包含每个$V_J$（$J\geq j_0$），由定理~\ref{thm:wavelet-dyadic-dense}在$L^2([0,1])$中稠密，所以它是完备的标准正交系。展开式与Parseval等式由第~\ref{chap:inner-product-spaces}章的正交展开得到。截到$j\leq J-1$的部分和是$f$在$V_{j_0}\oplus W_{j_0}\oplus\cdots\oplus W_{J-1}=V_J$上的正交投影，即$P_Jf$。
\end{proof}

每一层只有有限个系数，所以式~\eqref{eq:wavelet-haar-expansion}的外层是对$j$的级数，内层是有限和，不存在同层内求和次序的问题。按第~\ref{chap:inner-product-spaces}章的结论，标准正交展开的收敛也不依赖项的排列。

\paragraph{粗项不能省略}
每个Haar函数在$[0,1]$上的积分都是零。若去掉式~\eqref{eq:wavelet-haar-expansion}中的粗项，剩下的级数只能表示在每个第$j_0$层区间上平均为零的函数。常数函数$1$与所有$\psi_{j,k}$正交，没有粗项就完全无法表示。在区间上，粗项记录的是有限个区间上的平均水平，它的范数一般是正的，不能当作可以忽略的余项。第~\ref{sec:wavelet-real-line-expansion}小节会看到，在全实轴上情况不同。

\paragraph{有限变换的位置}
第~\ref{sec:wavelet-finite-haar}节的有限Haar变换，正是对$V_J$中的函数计算前$J$层系数：它不近似任何东西，只是在同一个有限维空间里换基。对一般的$f$，先用$P_Jf$（即区间平均）得到有限分辨率的表示，再作有限变换，得到的就是式~\eqref{eq:wavelet-haar-expansion}中$j<J$的全部系数。截断引起的误差是$\lVert f-P_Jf\rVert_2$，它由被丢掉的系数的平方和给出。

\paragraph{点态与一致收敛要另外核验}
定理~\ref{thm:wavelet-haar-expansion}的收敛是$L^2$意义下的，不保证每一点收敛，也不保证一致收敛。连续函数是一个可以直接验证的好情形：若$f$在$[0,1]$上连续，记$\omega(\delta)=\sup\{\lvert f(s)-f(t)\rvert:\lvert s-t\rvert\leq\delta\}$，则$x\in I_{j,k}$时
\[
\lvert P_jf(x)-f(x)\rvert=\Bigl\lvert2^j\int_{I_{j,k}}\bigl(f(t)-f(x)\bigr)\,dt\Bigr\rvert\leq\omega(2^{-j}).
\]
由一致连续性，$\omega(2^{-j})\to0$，所以$P_jf$一致收敛到$f$。

跳跃函数则不同。取$f=\mathbf{1}_{[1/3,1]}$。$1/3$不是二进点，所以每一层都有一个区间$I_{j,k}$包含$1/3$作为内点，$f$在其中既取$0$也取$1$，$P_jf$在该区间上取某个平均值$c_j\in(0,1)$。于是$\sup_x\lvert P_jf(x)-f(x)\rvert\geq\max(c_j,1-c_j)\geq1/2$，部分和对任何$j$都不一致逼近$f$，尽管$L^2$误差趋于零。对一般的$L^2$函数，部分和几乎处处收敛也成立，但证明需要Lebesgue微分定理，这里不给出证明。截断误差怎样随层数下降，要用函数的正则性定量刻画，留到第~\ref{chap:function-approximation}章。
